\section{有限温度 Feynman 规则}

Mahan 第 5 章把 Wick 定理落实为可逐图计算的规则。以下针对瞬时二体相互作用与 Matsubara 形式（$\hbar=k_B=1$）。

\subsection{相互作用绘景}

\begin{equation}
  \mathcal{G}(\tau,\tau')
  =-\frac{
  \bigl\langle
  T_\tau\,
  c(\tau)c^\dagger(\tau')
  S(\beta)
  \bigr\rangle_0
  }{
  \langle S(\beta)\rangle_0
  },
  \qquad
  S(\beta)
  =T_\tau\exp\Bigl(
  -\int_0^\beta\mathrm{d}\tau\,H_{\mathrm{int}}(\tau)
  \Bigr).
\end{equation}
分母消去不连通真空图（连锁集群）。$n$ 阶项含 $2n$ 个 $c$ 与 $2n$ 个 $c^\dagger$（每个 $H_{\mathrm{int}}$ 四个算符），由 Wick 完全收缩。

\subsection{图规则（动量--频率空间）}

对平移不变费米系统：
\begin{enumerate}
  \item 每条有向实线：$\mathcal{G}_0(\mathbf{k},\mathrm{i}\omega_n)=1/(\mathrm{i}\omega_n-\xi_{\mathbf{k}})$；
  \item 每条相互作用虚线：$v_{\mathbf{q}}$（或 $-U$ 等），在顶点守恒 $\mathbf{k}$、$\omega_n$、自旋；
  \item 对每个内部 $\mathbf{k},n$：$\dfrac{1}{\beta L^d}\sum_{\mathbf{k}n}$；
  \item 每个闭合费米圈：额外因子 $(-1)$，自旋求和常给 $2$；
  \item 对称因子：若图有 $n$ 个等价顶点，常含 $1/n!$ 与顶点数的抵消，最终只留图的对称群阶的倒数；
  \item 外腿不求和；自能图去掉两条外 $G_0$。
\end{enumerate}

\subsection{最低阶自能}

Hartree：
\begin{equation}
  \Sigma^{\mathrm{H}}
  =\frac{1}{\beta L^d}\sum_{\mathbf{q},m}
  v_{\mathbf{q}=0}\,
  e^{\mathrm{i}\omega_m 0^+}
  \mathcal{G}(\mathbf{p},\mathrm{i}\omega_m)
  \times(\text{自旋})
  =v_0 n_{\mathrm{tot}}.
\end{equation}
Fock：
\begin{equation}
  \Sigma^{\mathrm{F}}(\mathbf{k})
  =-\frac{1}{\beta L^d}\sum_{\mathbf{p},m}
  v_{\mathbf{k}-\mathbf{p}}
  e^{\mathrm{i}\omega_m 0^+}
  \mathcal{G}(\mathbf{p},\mathrm{i}\omega_m).
\end{equation}
用 \eqref{eq:G0_tau} 的 $0^+$ 约定完成频率求和后回到静态 HF。

\subsection{极化与屏蔽相互作用}

不可约极化（气泡）
\begin{equation}
  \Pi(\mathbf{q},\mathrm{i}\nu_m)
  =\frac{2}{\beta L^d}
  \sum_{\mathbf{k}n}
  \mathcal{G}(\mathbf{k},\mathrm{i}\omega_n)
  \mathcal{G}(\mathbf{k}+\mathbf{q},\mathrm{i}\omega_n+\mathrm{i}\nu_m)
  +\text{顶点修正}.
\end{equation}
RPA 取裸 $G_0$、无顶点修正，即 \eqref{eq:Pi_matsum}。屏蔽库仑线
\begin{equation}
  W(\mathbf{q},\mathrm{i}\nu_m)
  =\frac{v_{\mathbf{q}}}{1-v_{\mathbf{q}}\Pi(\mathbf{q},\mathrm{i}\nu_m)},
\end{equation}
与介电函数 $\varepsilon=1-v\Pi$ 一致。GW 自能 $\Sigma=\mathrm{i}G W$（实频）在 Matsubara 下为
\begin{equation}
  \Sigma(\mathbf{k},\mathrm{i}\omega_n)
  =-\frac{1}{\beta L^d}
  \sum_{\mathbf{q}m}
  W(\mathbf{q},\mathrm{i}\nu_m)
  \mathcal{G}(\mathbf{k}-\mathbf{q},\mathrm{i}\omega_n-\mathrm{i}\nu_m).
\end{equation}

\begin{note}
Mahan 的图规则与零温 Fetter--Walecka 规则平行：把 $\int\mathrm{d}t/2\pi$ 换成 $\beta^{-1}\sum_n$，把 $G_0(\omega)$ 换成 $\mathcal{G}_0(\mathrm{i}\omega_n)$。掌握一次，零温/有限温可以互相翻译。
\end{note}
